%% kapa: exam=1 round=35 year=2024 period=1 form=A subject=real_estate subject_ko=부동산학원론 seq=39 number=119
%% source: Q-Net 감정평가사 기출문제 (cst00309) artlSeq=5213229 file=raw/1008/5213229/1교시 문제지 원본.pdf
%% conversion: kordoc 4.24.51 -> tools/gen1.py (auto); review=proofread-sonnet; flags=none
\begin{question}{119}
다음은 매장의 매출액이 손익분기점 매출액 이하이면 기본임대료만 지급하고, 손익분기점 매출액 초과이면 초과매출액에 대하여 일정 임대료율을 적용한 추가임대료를 기본임대료에 가산하여 임대료를 지급하는 비율임대차(percentage lease)방식의 임대차계약의 조건이다. 이 임대차계약에서 계약기간 동안 지급할 것으로 예상되는 임대료의 합계는? (단, 주어진 조건에 한함)
\begin{framedtext}
○ 계약기간: 1년(1월 ∼ 12월) \par
○ 매장 임대면적: 200 $\mathrm{m}^{2}$ \par
○ 임대면적당 기본임대료: 월 5만원/$\mathrm{m}^{2}$ \par
○ 손익분기점 매출액: 월 2,000만원 \par
○ 각 월별 예상매출액 \par
\hspace*{1em}- 1월 ∼ 7월: 8만원/$\mathrm{m}^{2}$ \par
\hspace*{1em}- 8월 ∼ 12월: 20만원/$\mathrm{m}^{2}$ \par
○ 손익분기점 초과시 초과매출액에 대한 임대료율: 10 \% \par
\end{framedtext}
\begin{choices}
  \item 11,000만원
  \item 11,500만원
  \item 12,000만원
  \item 12,500만원
  \item 13,000만원
\end{choices}
\end{question}
\answer{5}
